% =====================================================================
% 6. Code distribution and reproduction (MLLB, Maria Martinez)
% README, DECISIONS D5, D11, D14.
% =====================================================================
\section{Code distribution and reproduction}\label{sec:code}

\subsection{Repository}

All code, logs, tables and figures are in the public repository
\url{https://github.com/orianafe16/ptbxl-DL}. The raw data and the cache are not committed: they are
downloaded from PhysioNet, checked against the published checksums and rebuilt. Checkpoints are not
committed either; predictions and logs are.

\begin{table}[htbp]
\centering
\caption{Structure of the repository and the files of my blocks.}
\label{tab:repo}
\small
\begin{tabularx}{\textwidth}{>{\raggedright\arraybackslash}p{4.2cm} >{\raggedright\arraybackslash}X}
\toprule
\textbf{Path} & \textbf{Content} \\
\midrule
\file{src/config.py}, \file{src/data.py}, \file{src/prepare_data.py}, \file{src/dataset.py} & Paths and settings; integrity check, labels and cache; split, normalisation, \texttt{Dataset} and \texttt{DataLoader}s (mine) \\
\file{src/metrics.py}, \file{src/tasks.py}, \file{src/train.py}, \file{src/runlog.py}, \file{src/compare.py} & Metrics and bootstrap, tasks, training loop, run logs, paired comparison (evaluation mine, loop shared) \\
\file{src/models/rnn.py}, \file{src/rnn_experiments.py} & LSTM, GRU and attention pooling; ablations, final models, report (mine) \\
\file{src/models/attention.py}, \file{src/models/transformer.py}, \file{src/transformer_experiments.py} & Hand-written attention, encoder blocks, Transformer classifier; ablations, final model, report (mine) \\
\file{src/models/dense.py}, \file{src/models/cnn.py}, generative and learning-strategy modules & Oriana's blocks and the shared area-2 experiments \\
\file{tests/} & Automated tests (pytest) of every module \\
\file{results/logs/}, \file{results/predictions/} & One JSON log and one prediction file per run; one summary per block \\
\file{report/} & Tables and figures generated by the code, \file{runs.json}, \file{make_numbers.py}, \file{numbers.tex}, these reports \\
\file{docs/} & Decisions log, task plan, AI usage log, one note per block, literature quotes \\
\bottomrule
\end{tabularx}
\end{table}

\subsection{Reproducing my results}

From a clean copy, with the Python environment of \file{requirements.txt} and the PTB-XL files in
\file{data/raw/ptbxl/}:
\begin{lstlisting}[language=bash,numbers=none,caption={Commands to reproduce my blocks.}]
python -m src.prepare_data                  # check the files, build the cache
python -m pytest -q tests                   # tests
python -m src.rnn_experiments estimate      # time per epoch on this GPU
python -m src.rnn_experiments all --fast    # smoke test (not reported)
python -m src.rnn_experiments all           # ablations + final models + report
python -m src.transformer_experiments estimate
python -m src.transformer_experiments all --fast
python -m src.transformer_experiments all
python report/make_numbers.py               # report/numbers.tex
\end{lstlisting}
\texttt{run} and \texttt{final} skip what has already been run, so an interrupted block can be resumed. The
cost of each configuration on the laptop is in Table~\ref{tab:mllb-cost}. Reproduction has three levels.
(1) From a clean clone, without data or GPU, the test suite runs on synthetic data, and
\file{report/make_numbers.py} rebuilds \file{numbers.tex} from the committed logs, which is all the reports
need; we checked that it gives exactly the same macros. (2) With the PTB-XL files and the checkpoints,
which are kept outside git and backed up separately, the \texttt{report} step of each block rebuilds its
tables and figures; the interpretation figures (attention weights, input sensitivity) need the
checkpoints. (3) Retraining repeats the runs; on a GPU, results can differ slightly from the logged ones,
because PyTorch's fully deterministic mode is not enabled.

\subsection{Tests}

The tests run on a CPU with synthetic data, so they need neither the GPU nor the dataset (the tests on the
real cache are skipped when it is missing). For my blocks they check, among other things:
\begin{itemize}[leftmargin=1.2em,itemsep=1pt]
  \item the data: split by fold, exclusion of records without a superclass, training-only statistics,
        class counts against the PhysioNet table, no patient in two splits;
  \item the evaluation: metrics on toy cases with known answers, patients kept together in the bootstrap,
        identical resamples for the same seed, a refused second final run;
  \item the recurrent models: framing element by element, LSTM and GRU states, the bidirectional final
        state, causality, the attention weights and the bound of the notebook score, the input sensitivity;
  \item the attention models: equality with \texttt{nn.MultiheadAttention}, the causal mask, the order of
        Pre-LN and Post-LN, the patch embedding equal to a \texttt{Conv1d}, the parameter count, the order
        invariance without positions, and the single-layer model equal to an attention pooling;
  \item for both blocks, the selection rule and the whole pipeline on a small synthetic cache.
\end{itemize}

\subsection{Traceability of the final results}

Table~\ref{tab:comparison_final} is the final comparison of the project, generated by
\texttt{python -m src.comparison report} from the summaries of every block. Each value is a macro, and each
macro can be traced to one key of one log listed in \file{report/runs.json}; each log records the commit of
the code that produced it and whether that code had uncommitted changes. The runs of my blocks were made
at committed versions of the code (commits \texttt{cae71d5} for the recurrent block and \texttt{f11e501} for
the Transformer). The only values typed by hand are the published benchmark rows, stored in
\file{report/published.json} with file, page and quote.

% Shared float (common/make_floats.py). Table: tables/comparison_final.tex
\begin{table}[htbp]
\centering
\caption[Final comparison (test)]{Final comparison on the test fold. Mean $\pm$ SD over seeds; interval: 95\,\% patient bootstrap of the first seed. Published rows: macro AUC of the superclass task in Table~I of \cite{strodthoff2021deep}, obtained in a different setting (described in the text).}
\label{tab:comparison_final}
\small
\begin{adjustbox}{max width=\textwidth}
% Generated by src/comparison.py from results/logs/comparison_summary.json.
% Mean +- SD over seeds (one run: no SD); CI = 95% patient bootstrap of the first seed. Needs alias "comparison".
\begin{tabular}{lcrrccc}
\toprule
Model & Labels & Train.\ param. & Seeds & Test macro AUROC & First seed [95\% CI] & Test macro F1 \\
\midrule
\multicolumn{7}{l}{\textit{12 leads}} \\
\quad Trivial (lead mean and SD) & 100\% & \comparisonParamsTrivial{} & \comparisonNSeedsTrivial{} & \comparisonTestAurocMeanTrivial{} & \comparisonTestAurocFirstTrivial{} [\comparisonTestAurocLoFirstTrivial{}, \comparisonTestAurocHiFirstTrivial{}] & \comparisonTestFOneMeanTrivial{} \\
\quad Dense & 100\% & \comparisonParamsTwelveDense{} & \comparisonNSeedsTwelveDense{} & \comparisonTestAurocMeanTwelveDense{} $\pm$ \comparisonTestAurocSdTwelveDense{} & \comparisonTestAurocFirstTwelveDense{} [\comparisonTestAurocLoFirstTwelveDense{}, \comparisonTestAurocHiFirstTwelveDense{}] & \comparisonTestFOneMeanTwelveDense{} $\pm$ \comparisonTestFOneSdTwelveDense{} \\
\quad CNN & 100\% & \comparisonParamsTwelveCnn{} & \comparisonNSeedsTwelveCnn{} & \comparisonTestAurocMeanTwelveCnn{} $\pm$ \comparisonTestAurocSdTwelveCnn{} & \comparisonTestAurocFirstTwelveCnn{} [\comparisonTestAurocLoFirstTwelveCnn{}, \comparisonTestAurocHiFirstTwelveCnn{}] & \comparisonTestFOneMeanTwelveCnn{} $\pm$ \comparisonTestFOneSdTwelveCnn{} \\
\quad RNN & 100\% & \comparisonParamsTwelveRnn{} & \comparisonNSeedsTwelveRnn{} & \comparisonTestAurocMeanTwelveRnn{} $\pm$ \comparisonTestAurocSdTwelveRnn{} & \comparisonTestAurocFirstTwelveRnn{} [\comparisonTestAurocLoFirstTwelveRnn{}, \comparisonTestAurocHiFirstTwelveRnn{}] & \comparisonTestFOneMeanTwelveRnn{} $\pm$ \comparisonTestFOneSdTwelveRnn{} \\
\quad RNN + attention & 100\% & \comparisonParamsTwelveRnnattn{} & \comparisonNSeedsTwelveRnnattn{} & \comparisonTestAurocMeanTwelveRnnattn{} $\pm$ \comparisonTestAurocSdTwelveRnnattn{} & \comparisonTestAurocFirstTwelveRnnattn{} [\comparisonTestAurocLoFirstTwelveRnnattn{}, \comparisonTestAurocHiFirstTwelveRnnattn{}] & \comparisonTestFOneMeanTwelveRnnattn{} $\pm$ \comparisonTestFOneSdTwelveRnnattn{} \\
\quad Transformer & 100\% & \comparisonParamsTwelveTransformer{} & \comparisonNSeedsTwelveTransformer{} & \comparisonTestAurocMeanTwelveTransformer{} $\pm$ \comparisonTestAurocSdTwelveTransformer{} & \comparisonTestAurocFirstTwelveTransformer{} [\comparisonTestAurocLoFirstTwelveTransformer{}, \comparisonTestAurocHiFirstTwelveTransformer{}] & \comparisonTestFOneMeanTwelveTransformer{} $\pm$ \comparisonTestFOneSdTwelveTransformer{} \\
\addlinespace
\multicolumn{7}{l}{\textit{Lead II only (T14)}} \\
\quad Dense & 100\% & \comparisonParamsOneleadDense{} & \comparisonNSeedsOneleadDense{} & \comparisonTestAurocMeanOneleadDense{} $\pm$ \comparisonTestAurocSdOneleadDense{} & \comparisonTestAurocFirstOneleadDense{} [\comparisonTestAurocLoFirstOneleadDense{}, \comparisonTestAurocHiFirstOneleadDense{}] & \comparisonTestFOneMeanOneleadDense{} $\pm$ \comparisonTestFOneSdOneleadDense{} \\
\quad CNN & 100\% & \comparisonParamsOneleadCnn{} & \comparisonNSeedsOneleadCnn{} & \comparisonTestAurocMeanOneleadCnn{} $\pm$ \comparisonTestAurocSdOneleadCnn{} & \comparisonTestAurocFirstOneleadCnn{} [\comparisonTestAurocLoFirstOneleadCnn{}, \comparisonTestAurocHiFirstOneleadCnn{}] & \comparisonTestFOneMeanOneleadCnn{} $\pm$ \comparisonTestFOneSdOneleadCnn{} \\
\quad RNN & 100\% & \comparisonParamsOneleadRnn{} & \comparisonNSeedsOneleadRnn{} & \comparisonTestAurocMeanOneleadRnn{} $\pm$ \comparisonTestAurocSdOneleadRnn{} & \comparisonTestAurocFirstOneleadRnn{} [\comparisonTestAurocLoFirstOneleadRnn{}, \comparisonTestAurocHiFirstOneleadRnn{}] & \comparisonTestFOneMeanOneleadRnn{} $\pm$ \comparisonTestFOneSdOneleadRnn{} \\
\quad RNN + attention & 100\% & \comparisonParamsOneleadRnnattn{} & \comparisonNSeedsOneleadRnnattn{} & \comparisonTestAurocMeanOneleadRnnattn{} $\pm$ \comparisonTestAurocSdOneleadRnnattn{} & \comparisonTestAurocFirstOneleadRnnattn{} [\comparisonTestAurocLoFirstOneleadRnnattn{}, \comparisonTestAurocHiFirstOneleadRnnattn{}] & \comparisonTestFOneMeanOneleadRnnattn{} $\pm$ \comparisonTestFOneSdOneleadRnnattn{} \\
\quad Transformer & 100\% & \comparisonParamsOneleadTransformer{} & \comparisonNSeedsOneleadTransformer{} & \comparisonTestAurocMeanOneleadTransformer{} $\pm$ \comparisonTestAurocSdOneleadTransformer{} & \comparisonTestAurocFirstOneleadTransformer{} [\comparisonTestAurocLoFirstOneleadTransformer{}, \comparisonTestAurocHiFirstOneleadTransformer{}] & \comparisonTestFOneMeanOneleadTransformer{} $\pm$ \comparisonTestFOneSdOneleadTransformer{} \\
\addlinespace
\multicolumn{7}{l}{\textit{Lead II, ECG-GPT transfer (T13)}} \\
\quad ECG-GPT, from scratch & 100\% & \comparisonParamsScratchFOneZeroZero{} & \comparisonNSeedsScratchFOneZeroZero{} & \comparisonTestAurocMeanScratchFOneZeroZero{} $\pm$ \comparisonTestAurocSdScratchFOneZeroZero{} & \comparisonTestAurocFirstScratchFOneZeroZero{} [\comparisonTestAurocLoFirstScratchFOneZeroZero{}, \comparisonTestAurocHiFirstScratchFOneZeroZero{}] & \comparisonTestFOneMeanScratchFOneZeroZero{} $\pm$ \comparisonTestFOneSdScratchFOneZeroZero{} \\
\quad ECG-GPT, linear probe (frozen) & 100\% & \comparisonParamsProbeFOneZeroZero{} & \comparisonNSeedsProbeFOneZeroZero{} & \comparisonTestAurocMeanProbeFOneZeroZero{} $\pm$ \comparisonTestAurocSdProbeFOneZeroZero{} & \comparisonTestAurocFirstProbeFOneZeroZero{} [\comparisonTestAurocLoFirstProbeFOneZeroZero{}, \comparisonTestAurocHiFirstProbeFOneZeroZero{}] & \comparisonTestFOneMeanProbeFOneZeroZero{} $\pm$ \comparisonTestFOneSdProbeFOneZeroZero{} \\
\quad ECG-GPT, fine-tuning & 100\% & \comparisonParamsFinetuneFOneZeroZero{} & \comparisonNSeedsFinetuneFOneZeroZero{} & \comparisonTestAurocMeanFinetuneFOneZeroZero{} $\pm$ \comparisonTestAurocSdFinetuneFOneZeroZero{} & \comparisonTestAurocFirstFinetuneFOneZeroZero{} [\comparisonTestAurocLoFirstFinetuneFOneZeroZero{}, \comparisonTestAurocHiFirstFinetuneFOneZeroZero{}] & \comparisonTestFOneMeanFinetuneFOneZeroZero{} $\pm$ \comparisonTestFOneSdFinetuneFOneZeroZero{} \\
\quad ECG-GPT, from scratch & 10\% & \comparisonParamsScratchFOneZero{} & \comparisonNSeedsScratchFOneZero{} & \comparisonTestAurocMeanScratchFOneZero{} $\pm$ \comparisonTestAurocSdScratchFOneZero{} & \comparisonTestAurocFirstScratchFOneZero{} [\comparisonTestAurocLoFirstScratchFOneZero{}, \comparisonTestAurocHiFirstScratchFOneZero{}] & \comparisonTestFOneMeanScratchFOneZero{} $\pm$ \comparisonTestFOneSdScratchFOneZero{} \\
\quad ECG-GPT, linear probe (frozen) & 10\% & \comparisonParamsProbeFOneZero{} & \comparisonNSeedsProbeFOneZero{} & \comparisonTestAurocMeanProbeFOneZero{} $\pm$ \comparisonTestAurocSdProbeFOneZero{} & \comparisonTestAurocFirstProbeFOneZero{} [\comparisonTestAurocLoFirstProbeFOneZero{}, \comparisonTestAurocHiFirstProbeFOneZero{}] & \comparisonTestFOneMeanProbeFOneZero{} $\pm$ \comparisonTestFOneSdProbeFOneZero{} \\
\quad ECG-GPT, fine-tuning & 10\% & \comparisonParamsFinetuneFOneZero{} & \comparisonNSeedsFinetuneFOneZero{} & \comparisonTestAurocMeanFinetuneFOneZero{} $\pm$ \comparisonTestAurocSdFinetuneFOneZero{} & \comparisonTestAurocFirstFinetuneFOneZero{} [\comparisonTestAurocLoFirstFinetuneFOneZero{}, \comparisonTestAurocHiFirstFinetuneFOneZero{}] & \comparisonTestFOneMeanFinetuneFOneZero{} $\pm$ \comparisonTestFOneSdFinetuneFOneZero{} \\
\addlinespace
\multicolumn{7}{l}{\textit{Lead II, zero-shot (T24)}} \\
\quad CLIP ViT-B/32, zero-shot & 0\% & --- & \comparisonNSeedsZeroshot{} & \comparisonTestAurocMeanZeroshot{} & \comparisonTestAurocFirstZeroshot{} [\comparisonTestAurocLoFirstZeroshot{}, \comparisonTestAurocHiFirstZeroshot{}] & \comparisonTestFOneMeanZeroshot{} \\
\addlinespace
\multicolumn{7}{l}{\textit{Published (quoted, different settings: see text)}} \\
\quad resnet1d\_wang, 12 leads~\cite{strodthoff2021deep} & & & & \comparisonPublishedStrodthoffResnetOneDWang{} & & \\
\quad xresnet1d101, 12 leads~\cite{strodthoff2021deep} & & & & \comparisonPublishedStrodthoffXresnetOneDOneZeroOne{} & & \\
\quad lstm, 12 leads~\cite{strodthoff2021deep} & & & & \comparisonPublishedStrodthoffLstm{} & & \\
\quad fcn\_wang, 12 leads~\cite{strodthoff2021deep} & & & & \comparisonPublishedStrodthoffFcnWang{} & & \\
\quad inception1d, 12 leads~\cite{strodthoff2021deep} & & & & \comparisonPublishedStrodthoffInceptionOneD{} & & \\
\quad lstm\_bidir, 12 leads~\cite{strodthoff2021deep} & & & & \comparisonPublishedStrodthoffLstmBidir{} & & \\
\quad Wavelet+NN, 12 leads~\cite{strodthoff2021deep} & & & & \comparisonPublishedStrodthoffWaveletNn{} & & \\
\quad ensemble, 12 leads~\cite{strodthoff2021deep} & & & & \comparisonPublishedStrodthoffEnsemble{} & & \\
\bottomrule
\end{tabular}

\end{adjustbox}
\end{table}

